\chapter{Section C: Problem Solving}\label{ch:problems}

There are 57 problems in five groups. Every problem is of Easy or Medium difficulty and
is written in the style of an online judge: you are given a structure and asked to
implement one operation on it. Each statement shows the function to complete (the
starter code), the constraints, the permitted operations, the required complexity, an
optional bonus for doing better, and sample tests. The code is plain C/C++: no global
variables, no \cc{static} and no \cc{const} are needed anywhere, and the solutions in
Chapter~\ref{ch:solproblems} follow the same rule.

\para{Interfaces assumed throughout} Array list: \cc{items[]} and \cc{size}. SLL:
Listing~\ref{lst:sll} (\cc{snode}, \cc{sll} with \cc{head}, \cc{tail}, \cc{size}). DLL:
Listing~\ref{lst:dll} (\cc{node}, \cc{list}). Stack and queue: Listings~\ref{lst:stack}
and~\ref{lst:queue} (\cc{stInit}, \cc{stPush}, \cc{stPop}, \cc{stTop}, \cc{stEmpty},
\cc{stSize}; \cc{qInit}, \cc{qEnq}, \cc{qDeq}, \cc{qFront}, \cc{qEmpty}, \cc{qSize}).
Deque: Listing~\ref{lst:deque} (\cc{dqInit}, \cc{dqPushFront}, \cc{dqPushBack},
\cc{dqPopFront}, \cc{dqPopBack}, \cc{dqFront}, \cc{dqBack}, \cc{dqEmpty}, \cc{dqSize}).
Every pop, top, front and back returns 1 on success and 0 when the structure is empty,
and passes the item through an \cc{int *}. When a problem names a function, use exactly
that signature; when it forbids an operation, even a correct use of it forfeits credit.

\section{Arrays and array lists}\label{sec:probarrays}

\problem{1}{String compression (run-length encoding)}{Easy}
A \emph{run} is a maximal block of equal consecutive characters. The run-length encoding
of a string writes, for each run, the character followed by the length of the run. Given
a string \cc{s} of letters, implement \cc{rleEncode}: it writes the encoding into
\cc{out} and returns the length of the encoding. The count is always written, even when
it is 1, and it may have several digits.

Then implement \cc{compress}, which does the same inside the array \cc{a} of \cc{n}
characters, but a run of length 1 is written as the bare character. It returns the new
length. No second array.
\begin{lstlisting}[style=starter]
int rleEncode(char *s, char *out);   // out must end with '\0'
int compress(char *a, int n);        // in place, returns the new length
\end{lstlisting}
\permitted Array indexing only. No \cc{strlen}, \cc{strcat}, \cc{sprintf} or \cc{itoa}:
write the digits of a count yourself.
\constraints $1\le n\le10^5$; the string contains only letters.
\required $O(n)$ time; \cc{compress} uses $O(1)$ extra space.
\bonus Implement \cc{rleDecode(char *e, char *out)}.
\begin{tests}
aaabccdddd & a3b1c2d4\\
wwwwwwwwwwwwbbb & w12b3\\
compress: aabcccccaaa & a2bc5a3 (length 7)\\
compress: abbbbbbbbbbbb & ab12 (length 4)\\
decode: x3y2z1 & xxxyyz\\
\end{tests}

\problem{2}{Matrix transpose}{Easy}
Given an $n\times m$ matrix \cc{a} ($n$ rows, $m$ columns), implement \cc{transpose},
which stores the $m\times n$ transpose in \cc{t}, so that $t_{ji}=a_{ij}$. Then implement
\cc{transposeInPlace} for a square $n\times n$ matrix.
\begin{lstlisting}[style=starter]
void transpose(int a[][10], int n, int m, int t[][10]);
void transposeInPlace(int a[][10], int n);
\end{lstlisting}
\permitted Array indexing only.
\constraints $1\le n,m\le10$.
\required $O(nm)$ time; \cc{transposeInPlace} uses $O(1)$ extra space.
\begin{tests}
{}[[1,2,3],[4,5,6]] & [[1,4],[2,5],[3,6]]\\
{[[1,2],[3,4]] (in place)} & [[1,3],[2,4]]\\
{}[[7]] & [[7]]\\
\end{tests}

\problem{3}{The missing number}{Easy}
Given an array \cc{a} of $m$ distinct integers in increasing order, taken from
$\{1,2,\dots,m+1\}$, exactly one number of that range is missing. Return it.
\begin{lstlisting}[style=starter]
int missingNumber(int a[], int m);
\end{lstlisting}
\permitted Array reads only. No sorting, no copy of the array, no counting array.
\constraints $1\le m\le10^4$.
\required $O(m)$ time and $O(1)$ extra space.
\bonus A second function \cc{missingNumberFast} in $O(\log m)$ time, using the fact
that the array is sorted.
\begin{tests}
1 2 4 5 6 7 & 3\\
1 2 3 4 5 6 8 9 10 & 7\\
2 & 1\\
1 & 2\\
\end{tests}

\problem{4}{A sorted array list}{Medium}
An array list \cc{slist} is kept in non-decreasing order at all times. Implement the
following operations so that the order is preserved. Use binary search wherever the
position of a value is needed. Write the helpers \cc{lowerBound} (index of the first
item $\ge x$) and \cc{upperBound} (index of the first item $>x$) first.
(i)~\cc{insert}: insert $x$ after the equal items; return 0 if the list is full, else 1.
(ii)~\cc{erase}: remove one occurrence of $x$; return 1 if found, else 0.
(iii)~\cc{search}: the index of the first occurrence of $x$, or $-1$.
(iv)~\cc{count}: the number of occurrences of $x$.
\begin{lstlisting}[style=starter]
typedef struct { int items[100]; int size; } slist;
int lowerBound(slist *l, int x);
int upperBound(slist *l, int x);
int insert(slist *l, int x);
int erase(slist *l, int x);
int search(slist *l, int x);
int count(slist *l, int x);
\end{lstlisting}
\permitted Array access to \cc{items} and \cc{size}. No linear search for the position of
a value.
\required Finding a position, \cc{search} and \cc{count}: $O(\log n)$. \cc{insert},
\cc{erase}: $O(\log n)$ comparisons and $O(n)$ moves.
\begin{tests}
list 2 4 4 4 7 9: insert 5 & 2 4 4 4 5 7 9\\
list 2 4 4 4 7 9: erase 4 & 2 4 4 7 9\\
list 2 4 4 4 7 9: search 4, search 5 & 1, -1\\
list 2 4 4 4 7 9: count 4, count 9, count 3 & 3, 1, 0\\
\end{tests}

\problem{5}{Rotate an array}{Easy}
Given an array \cc{a} of $n$ integers, rotate it to the right by $k$ positions, in place.
For example $[1,2,3,4,5,6,7]$ with $k=3$ becomes $[5,6,7,1,2,3,4]$. The value $k$ may be
larger than $n$.
\begin{lstlisting}[style=starter]
void rotate(int a[], int n, int k);
\end{lstlisting}
\permitted Array indexing and a scalar temporary. No second array.
\constraints $0\le n\le10^5$, $0\le k\le10^9$.
\required $O(n)$ time (shifting by one place $k$ times costs $O(nk)$ and earns partial
credit).
\bonus $O(1)$ extra space with three reversals.
\begin{tests}
{}[1,2,3,4,5,6,7], k=3 & [5,6,7,1,2,3,4]\\
{}[1,2], k=5 & [2,1]\\
{}[1,2,3,4,5,6,7,8,9], k=4 & [6,7,8,9,1,2,3,4,5]\\
\end{tests}

\problem{6}{Merge two sorted arrays}{Easy}
Given sorted (non-decreasing) arrays \cc{a} of length $n$ and \cc{b} of length $m$, fill
\cc{c} with all $n+m$ items in sorted order and return $n+m$. Among equal items, those of
\cc{a} come first.
\begin{lstlisting}[style=starter]
int mergeArrays(int a[], int n, int b[], int m, int c[]);
\end{lstlisting}
\permitted Array indexing only. No sorting routine.
\constraints $0\le n,m\le10^4$.
\required $O(n+m)$ time.
\begin{tests}
{}[1,3,5,7], [2,3,6] & [1,2,3,3,5,6,7]\\
{}[], [4,5] & [4,5]\\
{}[1,2], [] & [1,2]\\
\end{tests}

\problem{7}{Remove duplicates from a sorted array}{Easy}
Given a sorted array \cc{a} of $n$ integers, keep one copy of every value, in place, in
the first positions. Return the new length. The cells after the new length may hold
anything.
\begin{lstlisting}[style=starter]
int dedupSorted(int a[], int n);
\end{lstlisting}
\permitted Array indexing only.
\constraints $0\le n\le10^5$.
\required $O(n)$ time and $O(1)$ extra space.
\begin{tests}
{}[1,1,2,2,2,3,5,5] & length 4: [1,2,3,5]\\
{}[7] & length 1: [7]\\
{}[] & length 0\\
\end{tests}

\problem{8}{Two-sum in a sorted array}{Easy}
Given a sorted array \cc{a} of $n$ integers and a \cc{target}, find indices $i<j$ with
$a[i]+a[j]=\mathit{target}$. Store them in \cc{*pi} and \cc{*pj} and return 1, or return
0 if there is no such pair.
\begin{lstlisting}[style=starter]
int twoSum(int a[], int n, int target, int *pi, int *pj);
\end{lstlisting}
\permitted Array indexing only (no hash table).
\constraints $2\le n\le10^5$; values fit in \cc{int} and so do the sums.
\required $O(n^2)$ earns partial credit; $O(n)$ time and $O(1)$ space earns full credit.
\begin{tests}
{}[1,3,4,6,8,11], target 10 & (2,3)\\
{}[1,3,4,6,8,11], target 20 & none\\
{}[2,2], target 4 & (0,1)\\
\end{tests}

\problem{9}{Maximum subarray sum}{Medium}
Given an array \cc{a} of $n\ge1$ integers (positive, negative or zero), return the largest
sum of a non-empty contiguous subarray.
\begin{lstlisting}[style=starter]
int maxSubarray(int a[], int n);
\end{lstlisting}
\permitted Array reads only.
\constraints $1\le n\le10^5$; $|a[i]|\le10^4$.
\required $O(n^2)$ earns partial credit (try all subarrays with a running sum); $O(n)$
time and $O(1)$ space earns full credit.
\begin{tests}
{}[-2,1,-3,4,-1,2,1,-5,4] & 6 (subarray [4,-1,2,1])\\
{}[-3,-1,-2] & -1\\
{}[5] & 5\\
\end{tests}

\problem{10}{Move zeroes}{Easy}
Given an array \cc{a} of $n$ integers, move all zeros to the end, keeping the order of the
non-zero items, in place.
\begin{lstlisting}[style=starter]
void moveZeroes(int a[], int n);
\end{lstlisting}
\permitted Array indexing only.
\constraints $0\le n\le10^5$.
\required $O(n)$ time, $O(1)$ extra space.
\begin{tests}
{}[0,1,0,3,12] & [1,3,12,0,0]\\
{}[0,0] & [0,0]\\
{}[4,5] & [4,5]\\
\end{tests}

\problem{11}{Majority element}{Easy}
Given an array \cc{a} of $n$ integers, return the majority element, the element that
appears more than $\lfloor n/2\rfloor$ times. It is guaranteed to exist.
\begin{lstlisting}[style=starter]
int majority(int a[], int n);
\end{lstlisting}
\permitted Array reads only. No sorting, no counting array indexed by value.
\constraints $1\le n\le10^5$.
\required $O(n)$ time.
\bonus $O(1)$ extra space (Boyer--Moore voting).
\begin{tests}
{}[2,2,1,1,2,2,3,2,2] & 2\\
{}[3,3,4] & 3\\
{}[7] & 7\\
\end{tests}

\problem{12}{Rotate a matrix by 90 degrees}{Medium}
Given an $n\times n$ matrix \cc{a}, rotate it by $90^\circ$ clockwise, in place.
\begin{lstlisting}[style=starter]
void rotate90(int a[][10], int n);
\end{lstlisting}
\permitted Array indexing only.
\constraints $1\le n\le10$.
\required $O(n^2)$ time, $O(1)$ extra space.
\begin{tests}
{}[[1,2],[3,4]] & [[3,1],[4,2]]\\
{}[[1,2,3],[4,5,6],[7,8,9]] & [[7,4,1],[8,5,2],[9,6,3]]\\
\end{tests}

\section{Linked lists}\label{sec:problinked}

\problem{13}{Implementation: eraseItems (DLL)}{Medium}
\emph{This is one of the implementation problems of the departmental style, reproduced
with its sample input and output.}

Implement the operation \cc{void eraseItems(list *l, int x)} for a doubly linked list.
Without calling \cc{search} and \cc{erase}, remove all items in \cc{l} equal to \cc{x}.
The type \cc{list} (Listing~\ref{lst:dll}) has the fields \cc{head}, \cc{tail},
\cc{size}; nodes have \cc{item}, \cc{next}, \cc{prev}.

The test program reads $n$, then $n$ items, then $x$. It prints the list before, calls
your function, prints the list after, then \cc{New size:} and the list from tail to head.
When the list becomes empty it also prints a check that \cc{head} and \cc{tail} are both
\cc{NULL}. (The template code around your function is reconstructed from the sample
output; here the test cases are separated by a line \cc{----}.)
\begin{lstlisting}[style=starter]
void eraseItems(list *l, int x);
\end{lstlisting}
\permitted Direct access to \cc{head}, \cc{tail}, \cc{size}, \cc{next}, \cc{prev}. Not
\cc{search}, \cc{erase}, \cc{remove} or any helper that finds or deletes a node for you.
\required $O(n)$ time in one pass; $O(1)$ extra space; \cc{size}, \cc{head}, \cc{tail}
and every \cc{prev} link correct afterwards; removed nodes freed.
\bonus Explain why a solution that restarts from the head after every removal is
$\Theta(n^2)$ on an input such as \cc{1 1 1 1 1 3 3 3 3 3} with $x=3$, although
it is fast on \cc{3 3 3 3 3}.

\par\smallskip\noindent\textsc{\small Input} (the test cases one after another)
\VerbatimInput[fontsize=\footnotesize,frame=single,rulecolor=\color{rulegray}]{\codepath/notes/p13_template.in}
\noindent\textsc{\small Expected output}
\VerbatimInput[fontsize=\footnotesize,frame=single,rulecolor=\color{rulegray}]{\codepath/out/p13_template.txt}

\problem{14}{Implementation: appendList (copying a list)}{Medium}
Implement the function \cc{void appendList(list *l, list m)}. The idea is simply to
append the items of \cc{m} to \cc{l}, without calling \cc{append}. Take note of the edge
cases. Make sure that when \cc{m} is destroyed, the items appended to the list pointed to
by \cc{l} are no longer affected. The type \cc{list} (Listing~\ref{lst:dll}) has the
fields \cc{head}, \cc{tail}, \cc{size}; nodes have \cc{item}, \cc{next}, \cc{prev}.

The test program reads $n$ and the items of \cc{l}, then $|m|$ and the items of \cc{m};
it prints \cc{list l:}, \cc{l}, \cc{list m:}, \cc{m}, calls the function, destroys
\cc{m}, and prints \cc{l}.
\begin{lstlisting}[style=starter]
void appendList(list *l, list m);
\end{lstlisting}
\permitted Direct access to \cc{head}, \cc{tail}, \cc{size}, \cc{next}, \cc{prev};
\cc{malloc}. Not \cc{append}.
\required $O(|m|)$ time; the nodes of \cc{l} must be new nodes, never nodes of \cc{m}.
\bonus The call \cc{appendList(\&x, x)} (a list appended to itself) must terminate and
double the list. Say why a loop until \cc{NULL} fails here and a counted loop does not.

\par\smallskip\noindent\textsc{\small Input} (the test cases one after another)
\VerbatimInput[fontsize=\footnotesize,frame=single,rulecolor=\color{rulegray}]{\codepath/notes/p14_template.in}
\noindent\textsc{\small Expected output}
\VerbatimInput[fontsize=\footnotesize,frame=single,rulecolor=\color{rulegray}]{\codepath/out/p14_template.txt}

\problem{15}{Implementation: recursive append with a double pointer}{Medium}
Assume the singly linked list has no tail pointer. Implement \cc{append(snode **h, int
x)}, which adds a node with item \cc{x} at the end of the list. Notice the double pointer:
it lets the function change the real pointer and not a copy. \cc{h} is the address of the
pointer that heads the list. Implement it recursively. In \cc{main} it is called as
\cc{append(\&l.head, item)}: since the head must change when the list is empty, what is
passed is the address of the head and not the head.
\begin{lstlisting}[style=starter]
void append(snode **h, int x);   // called as append(&l.head, item)
\end{lstlisting}
\permitted \cc{malloc}; the fields \cc{item} and \cc{next}. No loops in \cc{append}, no
extra helper that returns a new head.
\required $O(n)$ time, recursive.
\bonus Implement the same operation with a loop but still with a \cc{snode **} (one
variable, no special case for the empty list). State the recursion depth and when it is
dangerous.
\begin{tests}
5 / 1 2 3 4 5 & 1 2 3 4 5\\
1 / 109 & 109\\
10 / 10 9 8 7 6 5 4 3 2 1 & 10 9 8 7 6 5 4 3 2 1\\
\end{tests}

\problem{16}{Reverse a singly linked list}{Easy}
Given an SLL, reverse it in place by changing links: do not allocate nodes and do not
copy items. Implement the iterative \cc{reverseList}, which must also fix \cc{head} and
\cc{tail}.
\begin{lstlisting}[style=starter]
void reverseList(sll *l);          // iterative; fix head and tail
snode *reverseRec(snode *h);       // bonus: recursive, returns the new head
\end{lstlisting}
\permitted Pointer manipulation of \cc{next}, \cc{head}, \cc{tail}.
\constraints $0\le|l|\le10^5$.
\required $O(n)$ time, $O(1)$ extra space (iterative).
\bonus The recursive version \cc{reverseRec}; explain its stack use.
\begin{tests}
{}[1,2,3,4,5] & [5,4,3,2,1]\\
{}[] & []\\
{}[7] & [7]\\
\end{tests}

\problem{17}{Middle node and $k$-th node from the end}{Easy}
Given an SLL, implement (a)~\cc{middleItem}, which returns the middle item (the second
middle when $n$ is even; the list is non-empty), and (b)~\cc{kthFromEnd}, which stores
the $k$-th item from the end in \cc{*out} and returns 1, or returns 0 if $k<1$ or $k>n$.
Both must work in one pass and without using \cc{size}.
\begin{lstlisting}[style=starter]
int middleItem(sll *l);                    // second middle if n is even
int kthFromEnd(sll *l, int k, int *out);   // k = 1 is the last item
\end{lstlisting}
\permitted Pointer traversal only. You may not read \cc{size}, count the nodes first, or
copy the list into an array.
\constraints $1\le n\le10^5$.
\required $O(n)$ time, one traversal, $O(1)$ extra space.
\bonus Return the \emph{first} middle for even $n$ by a change of one condition, and
explain the change.
\begin{tests}
middle [1,2,3,4] & 3\\
middle [1,2,3,4,5] & 3\\
kth [10,20,30,40,50], k=2 & 40\\
kth [10,20,30,40,50], k=5 / k=6 / k=0 & 10 / fails / fails\\
\end{tests}

\problem{18}{Splicing singly linked lists}{Medium}
Splicing moves nodes from one list to another by relinking, without allocating or
copying.
(i)~\cc{spliceAfter}: move all nodes of \cc{m} into \cc{l}, immediately after node
\cc{p} (\cc{p==NULL} means ``at the very front''). Afterwards \cc{m} is empty;
\cc{head}, \cc{tail} and both sizes must be correct.
(ii)~\cc{splice}: the same, but the place is an index: \cc{pos} nodes of \cc{l} stay in
front of the spliced part ($0\le\cc{pos}\le|l|$). Return 0 if \cc{pos} is invalid.
\begin{lstlisting}[style=starter]
void spliceAfter(sll *l, snode *p, sll *m);  // p == NULL: at the front
int splice(sll *l, int pos, sll *m);          // 0 if pos is invalid
\end{lstlisting}
\permitted Pointer relinking only. No \cc{malloc}, no \cc{free}, no item copying.
\required \cc{spliceAfter}: $O(1)$. \cc{splice}: $O(\cc{pos})$ to walk to the anchor,
then $O(1)$.
\bonus Show that, given a tail pointer, appending a whole list to another list costs
$O(1)$ whereas an array list needs $\Theta(|m|)$.
\begin{tests}
l=[1,2,5], m=[3,4], splice at pos 2 & l=[1,2,3,4,5] m=[] sizes 5 0\\
l=[1,2,5], m=[3,4], splice at pos 0 & l=[3,4,1,2,5]\\
l=[1,2,5], m=[3,4], splice at pos 3 & l=[1,2,5,3,4] (tail is 4)\\
m empty & l unchanged\\
\end{tests}

\problem{19}{Merge two sorted doubly linked lists}{Medium}
Two doubly linked lists \cc{a} and \cc{b} are sorted in non-decreasing order. Implement
\cc{mergeSorted} so that afterwards \cc{a} holds all the nodes in sorted order and
\cc{b} is empty. Relink the existing nodes (no allocation, no freeing); on equal items
the node of \cc{a} goes first; all \cc{prev} links, \cc{head}, \cc{tail} and \cc{size}
must be correct.
\begin{lstlisting}[style=starter]
void mergeSorted(list *a, list *b);
\end{lstlisting}
\permitted Pointer manipulation of \cc{next}, \cc{prev}, \cc{head}, \cc{tail}, \cc{size}.
\required $O(|a|+|b|)$ time, $O(1)$ extra space.
\bonus Attach the remainder of the non-empty list in $O(1)$ instead of node by node.
\begin{tests}
a=[1,4,7], b=[2,4,5,9] & a=[1,2,4,4,5,7,9] size 7; backward: 9 7 5 4 4 2 1; b.size=0\\
a=[], b=[3] & a=[3]\\
a=[1], b=[] & a=[1]\\
\end{tests}

\problem{20}{Palindromic doubly linked list}{Easy}
Given a doubly linked list, return 1 if its items read the same forwards and backwards,
else 0. The empty list and a one-element list are palindromes.
\begin{lstlisting}[style=starter]
int isPalindrome(list *l);
\end{lstlisting}
\permitted Pointer traversal. Not an auxiliary array or stack.
\constraints $0\le n\le10^5$.
\required $O(n)$ time, $O(1)$ extra space.
\bonus Stop after $\lfloor n/2\rfloor$ comparisons.
\begin{tests}
{}[1,2,3,2,1] & 1\\
{}[1,2,3,4] & 0\\
{}[] & 1\\
{}[7] & 1\\
{}[4,5,5,4] & 1\\
\end{tests}

\problem{21}{Reverse a doubly linked list}{Easy}
Given a doubly linked list, reverse it in place.
\begin{lstlisting}[style=starter]
void reverseDLL(list *l);
\end{lstlisting}
\permitted Pointer manipulation only.
\required $O(n)$ time, $O(1)$ space.
\bonus Do it by swapping \cc{next} and \cc{prev} in every node, without moving any item.
\begin{tests}
{}[1,2,3,4,5] & forward 5 4 3 2 1; backward 1 2 3 4 5\\
{}[] & []\\
\end{tests}

\problem{22}{Remove duplicates from a sorted doubly linked list}{Easy}
Given a doubly linked list sorted in non-decreasing order, remove nodes so that each
value appears once (keep the first node of each run). Free the removed nodes and keep
\cc{tail} and \cc{size} correct.
\begin{lstlisting}[style=starter]
void removeDuplicates(list *l);
\end{lstlisting}
\permitted Pointer manipulation only; no \cc{search} or \cc{erase}.
\required $O(n)$ time, $O(1)$ space.
\bonus Remove \emph{all} nodes whose value occurs more than once.
\begin{tests}
{}[1,1,2,3,3,3,5] & [1,2,3,5], size 4\\
{}[2,2,2] & [2]\\
{}[] & []\\
\end{tests}

\problem{23}{Insert into a sorted doubly linked list}{Easy}
Given a doubly linked list sorted in non-decreasing order, insert $x$ so that the list
stays sorted (after equal items). Handle the empty list, the new minimum and the new
maximum.
\begin{lstlisting}[style=starter]
void insertSorted(list *l, int x);
\end{lstlisting}
\permitted Allocation, \cc{next}/\cc{prev}, \cc{head}/\cc{tail}. No call to
\cc{append} except where it is the correct case.
\required $O(n)$ time.
\bonus $O(1)$ when $x$ is not smaller than the tail item.
\begin{tests}
{}[2,4,4,8], insert 1 & [1,2,4,4,8]\\
then insert 5 & [1,2,4,4,5,8]\\
then insert 9 & [1,2,4,4,5,8,9] (tail 9, size 7)\\
then insert 4 & [1,2,4,4,4,5,8,9]\\
{}[], insert 3 & [3]\\
\end{tests}

\problem{24}{An SLL with a middle pointer}{Medium}
A singly linked list \cc{mlist} keeps, besides \cc{head}, \cc{tail} and \cc{size}, a
pointer \cc{mid}. The invariant is
\begin{gather*}\cc{size}=0\Rightarrow\cc{head}=\cc{mid}=\cc{tail}=\cc{NULL},\\
\cc{size}>0\Rightarrow\cc{mid}\text{ points to the node at index }
\lfloor(\cc{size}-1)/2\rfloor\text{ (0-based).}\end{gather*}
Implement \cc{middle}, \cc{mPushBack}, \cc{mPushFront}, \cc{mPopFront} and \cc{mPopBack}
so that every operation restores the invariant.
\begin{lstlisting}[style=starter]
typedef struct mnode { int item; struct mnode *next; } mnode;
typedef struct { mnode *head, *mid, *tail; int size; } mlist;
void minit(mlist *l);
int middle(mlist *l);                 // O(1)
void mPushBack(mlist *l, int x);      // O(1)
int mPopFront(mlist *l, int *x);      // O(1); 0 if empty
void mPushFront(mlist *l, int x);     // O(1) or O(n), see below
int mPopBack(mlist *l, int *x);       // O(n)
\end{lstlisting}
\permitted Pointer manipulation. A node has only \cc{item} and \cc{next}.
\required \cc{middle}, \cc{mPushBack}, \cc{mPopFront}: $O(1)$. For \cc{mPushFront} and
\cc{mPopBack} state the worst case.
\bonus Explain why a doubly linked list would make all four end operations $O(1)$.
\begin{tests}
pushBack 10,20,\dots,70 (sizes 1..7) & mid after each: 10 10 20 20 30 30 40\\
10,20,30,40; pushFront 5; pushFront 1 & mid 20; then mid 10 (list 1 5 10 20 30 40)\\
then: popFront; popBack; pushBack 50; popFront & mid 20; 10; 20; 20 (list 10 20 30 50)\\
\end{tests}

\problem{25}{Rotate a linked list}{Medium}
Given an SLL and an integer $k\ge0$, rotate the list to the left by $k$ positions:
$[1,2,3,4,5]$ with $k=2$ becomes $[3,4,5,1,2]$. $k$ may exceed $n$.
\begin{lstlisting}[style=starter]
void rotateLeft(sll *l, int k);
\end{lstlisting}
\permitted Pointer manipulation; \cc{size} is available.
\constraints $0\le n\le10^5$, $0\le k\le10^9$.
\required $O(n)$ time.
\bonus No node is allocated or freed.
\begin{tests}
{}[1,2,3,4,5], k=2 & [3,4,5,1,2]\\
{}[1,2,3,4,5], k=7 & [3,4,5,1,2]\\
{}[1,2,3,4,5], k=5 & [1,2,3,4,5]\\
{}[], k=3 & []\\
\end{tests}

\problem{26}{Partition a linked list around a value}{Medium}
Given an SLL and a value $x$, rearrange the nodes so that all nodes with item $<x$
precede all nodes with item $\ge x$, keeping the original relative order inside each
group. Relink nodes; do not allocate.
\begin{lstlisting}[style=starter]
void partitionList(sll *l, int x);
\end{lstlisting}
\permitted Pointer manipulation only.
\required $O(n)$ time, one pass.
\bonus $O(1)$ extra space using two dummy nodes.
\begin{tests}
{}[1,4,3,2,5,2], x=3 & [1,2,2,4,3,5]\\
{}[5,6], x=3 & [5,6]\\
{}[1,2], x=3 & [1,2]\\
{}[], x=3 & []\\
\end{tests}

\problem{27}{Remove the $n$-th node from the end}{Medium}
Given an SLL, remove the $n$-th node from the end ($n=1$ is the last node) in one pass.
Return 1, or 0 if $n<1$ or $n>|l|$. Keep \cc{head}, \cc{tail} and \cc{size} correct and
free the node.
\begin{lstlisting}[style=starter]
int removeNthFromEnd(sll *l, int n);   // n = 1 is the last node
\end{lstlisting}
\permitted Pointer manipulation; do not use \cc{size} to compute the position.
\required $O(|l|)$ time, one pass.
\bonus Handle the removal of the head with a dummy node and no special case.
\begin{tests}
{}[1,2,3,4,5], n=2 & [1,2,3,5]\\
{}[1,2,3,5], n=4 & [2,3,5] (size 3)\\
{}[2,3,5], n=1 & [2,3] (tail 3)\\
{}[2,3], n=3 & fails (0), list unchanged\\
{}[7], n=1 & [] (head, tail NULL, size 0)\\
\end{tests}

\problem{28}{Swap nodes in pairs}{Medium}
Given an SLL, swap every two adjacent nodes: $[1,2,3,4]\to[2,1,4,3]$. A final unpaired
node stays where it is. Swap nodes, not items.
\begin{lstlisting}[style=starter]
void swapPairs(sll *l);
\end{lstlisting}
\permitted Pointer manipulation only.
\required $O(n)$ time, $O(1)$ space.
\bonus The same with recursion; also keep \cc{tail} correct.
\begin{tests}
{}[1,2,3,4] & [2,1,4,3]\\
{}[1,2,3,4,5] & [2,1,4,3,5] (tail 5)\\
{}[1] & [1]\\
\end{tests}

\problem{29}{Cycle detection}{Medium}
Given the head of a singly linked list whose last node may point back to an earlier node,
return 1 if the list contains a cycle and 0 otherwise. The list is not modified.
\begin{lstlisting}[style=starter]
int hasCycle(snode *head);        // 1 if the list has a cycle
snode *cycleStart(snode *head);   // bonus: first node of the cycle, or NULL
\end{lstlisting}
\permitted Pointer traversal. No visited-marker field, no auxiliary array or table.
\constraints $0\le n\le10^5$.
\required $O(n)$ time and $O(1)$ extra space.
\bonus Return the first node of the cycle (or \cc{NULL}), in $O(n)$ time and $O(1)$
space.
\begin{tests}
{}[1,2,3,4,5], last node points to the 3rd node & 1 (cycle starts at item 3)\\
{}[1,2,3,4,5], last node points to the head & 1 (starts at item 1)\\
{}[1,2,3,4,5], no link back & 0\\
{[] (head is NULL)} & 0\\
\end{tests}

\problem{30}{Group odd and even positions}{Medium}
Given an SLL, group all nodes at odd positions (1st, 3rd, \dots) first, followed by all
nodes at even positions, each group keeping its order: $[1,2,3,4,5]\to[1,3,5,2,4]$. The
positions are meant, not the values.
\begin{lstlisting}[style=starter]
void oddEvenList(sll *l);
\end{lstlisting}
\permitted Pointer manipulation only.
\required $O(n)$ time, $O(1)$ extra space.
\bonus Keep \cc{tail} correct.
\begin{tests}
{}[1,2,3,4,5] & [1,3,5,2,4] (tail 4)\\
{}[1,2,3,4] & [1,3,2,4]\\
{}[1,2] & [1,2]\\
\end{tests}

\problem{31}{Add two numbers stored in linked lists}{Medium}
Two non-negative integers are stored one digit per node, least significant digit first:
342 is $2\to4\to3$. Given the lists \cc{a} and \cc{b}, build in \cc{out} the list of their
sum, in the same form. The lists may have different lengths and there may be a final
carry.
\begin{lstlisting}[style=starter]
void addNumbers(sll *a, sll *b, sll *out);
\end{lstlisting}
\permitted \cc{sinit} and \cc{pushBack} on the result. No conversion to a machine
integer (the numbers may have hundreds of digits).
\required $O(\max(|a|,|b|))$ time.
\bonus The case in which the digits are stored most significant first.
\begin{tests}
{}[2,4,3] + [5,6,4] & [7,0,8]\\
{}[9,9] + [1] & [0,0,1]\\
{}[0] + [0] & [0]\\
\end{tests}

\section{Stacks}\label{sec:probstacks}

In this section the stack of Listing~\ref{lst:stack} holds \cc{int}s; characters are
pushed as their codes. ``Stack operations only'' means \cc{stInit}, \cc{stPush},
\cc{stPop}, \cc{stTop}, \cc{stEmpty}, \cc{stSize}: no access to \cc{items[]}.

\problem{32}{Valid brackets}{Easy}
Given a string \cc{s} containing only the characters \cc{( ) [ ] \{ \}}, return 1 if
every bracket is closed by the same kind of bracket and in the correct order, else 0.
\begin{lstlisting}[style=starter]
int isBalanced(char s[]);
\end{lstlisting}
\permitted Stack operations only.
\constraints $0\le|s|\le10^4$.
\required $O(n)$ time.
\begin{tests}
\{[()]\} & 1\\
([)] & 0\\
(( & 0\\
)) & 0\\
(empty string) & 1\\
()[]\{\} & 1\\
\end{tests}

\problem{33}{Infix to postfix and postfix evaluation}{Medium}
(a)~Given an infix expression \cc{in} made of single-digit operands, the operators
\cc{+ - * /} and parentheses (no spaces), write its postfix form into \cc{out}. The
operators \cc{*} and \cc{/} have higher precedence than \cc{+} and \cc{-}; operators of
equal precedence are left-associative.
(b)~Given a valid postfix expression \cc{p} of single digits and \cc{+ - * /}, return its
value (integer division, truncating toward zero).
\begin{lstlisting}[style=starter]
void toPostfix(char in[], char out[]);
int evalPostfix(char p[]);
\end{lstlisting}
\permitted Stack operations only.
\constraints $1\le|\mathit{in}|\le100$; the input is valid.
\required $O(n)$ time for both.
\bonus Explain what changes if \cc{\^} (exponent, right-associative) is added.
\begin{tests}
2+3*4 & 234*+ ; value 14\\
(2+3)*4 & 23+4* ; value 20\\
2*(3+4)-5 & 234+*5- ; value 9\\
8-3-2 & 83-2- ; value 3\\
9-8/(4-2) & 9842-/- ; value 5\\
\end{tests}

\problem{34}{A stack with constant-time minimum}{Easy}
Design a stack that supports \cc{push}, \cc{pop} and \cc{min} (the smallest item
currently on the stack), each in $O(1)$ time.
\begin{lstlisting}[style=starter]
typedef struct { stack all; stack mins; } minStack;
void msInit(minStack *m);
void msPush(minStack *m, int x);
int msPop(minStack *m, int *x);    // 0 if empty
int msMin(minStack *m, int *x);    // 0 if empty
\end{lstlisting}
\permitted Stack operations on the two member stacks.
\required Every operation $O(1)$.
\bonus Use less memory when many pushes are not new minima (push to \cc{mins} only when
$x\le$ the current minimum), and explain why the test must be $\le$, not $<$.
\begin{tests}
push 5 3 7 3 1; min & 1\\
pop; pop; min & 3 (stack 5 3 7)\\
push 2; min; pop; min & 2; 3\\
min on an empty stack & fails (0)\\
\end{tests}

\problem{35}{Next greater element}{Medium}
For every item \cc{a[i]} of an array, find the first item to its right that is
\emph{strictly} greater, or $-1$ if there is none.
\begin{lstlisting}[style=starter]
void nextGreater(int a[], int n, int out[]);
\end{lstlisting}
\permitted Array access to \cc{a} and \cc{out}; one stack.
\constraints $0\le n\le10^5$.
\required $O(n)$ time ($O(n^2)$ earns partial credit).
\bonus The \emph{circular} version: the search may wrap around to the beginning of the
array (\cc{[1,2,1]} $\to$ \cc{[2,-1,2]}).
\begin{tests}
{}[2,1,5,3] & [5,5,-1,-1]\\
{}[2,7,3,5,4,6,8] & [7,8,5,6,6,8,-1]\\
{}[5,4,3] & [-1,-1,-1]\\
{}[1,1] & [-1,-1]\\
\end{tests}

\problem{36}{Daily temperatures}{Medium}
Given the daily temperatures \cc{t[0..n-1]}, compute \cc{wait[i]}, the number of days
after day $i$ until a \emph{strictly} warmer day, or 0 if no warmer day follows.
\begin{lstlisting}[style=starter]
void dailyTemperatures(int t[], int n, int wait[]);
\end{lstlisting}
\permitted Array access to \cc{t} and \cc{wait}; one stack.
\constraints $1\le n\le10^5$; $30\le t[i]\le100$.
\required $O(n)$ time.
\begin{tests}
{}[73,74,75,71,69,72,76,73] & [1,1,4,2,1,1,0,0]\\
{}[30,40,50,60] & [1,1,1,0]\\
{}[30,60,90] & [1,1,0]\\
{}[50,50] & [0,0]\\
\end{tests}

\problem{37}{Online stock span}{Medium}
Prices of a stock arrive one per day. The \emph{span} of today's price is the number of
consecutive days, ending today and going backwards, on which the price was $\le$ today's
price. Implement \cc{ssNext}, which receives today's price and returns its span. The
prices are given one at a time: you cannot look at future prices, and you should not
store the whole history.
\begin{lstlisting}[style=starter]
typedef struct { stack prices; stack spans; } stockSpanner;
void ssInit(stockSpanner *s);
int ssNext(stockSpanner *s, int price);
\end{lstlisting}
\permitted Stack operations on the two member stacks.
\required $O(1)$ amortised per call.
\begin{tests}
100 80 60 70 60 75 85 & 1 1 1 2 1 4 6\\
10 10 10 & 1 2 3\\
5 4 3 2 1 & 1 1 1 1 1\\
\end{tests}

\problem{38}{Asteroid collision}{Medium}
Each item of \cc{a} is an asteroid in a row: its absolute value is its size, its sign its
direction ($>0$ moves right, $<0$ moves left); all move at the same speed. When two
meet, the smaller one explodes; if both are the same size, both explode. Two asteroids
moving in the same direction never meet. Write the survivors, left to right, into
\cc{out} and return how many there are.
\begin{lstlisting}[style=starter]
int asteroidCollision(int a[], int n, int out[]);
\end{lstlisting}
\permitted Array access to \cc{a} and \cc{out}; one stack.
\constraints $0\le n\le10^4$; $a[i]\ne0$.
\required $O(n)$ time.
\begin{tests}
{}[5,10,-5] & [5,10]\\
{}[8,-8] & []\\
{}[10,2,-5] & [10]\\
{}[-2,-1,1,2] & [-2,-1,1,2]\\
\end{tests}

\problem{39}{Decode string}{Medium}
A string is encoded by the rule \cc{k[part]}: the part inside the brackets is repeated
exactly $k$ times ($k\ge1$, possibly several digits). Encodings may be nested. Write the
decoded string into \cc{out}.
\begin{lstlisting}[style=starter]
void decodeString(char s[], char out[]);
\end{lstlisting}
\permitted Character access to \cc{s} and \cc{out}; stacks.
\constraints The input is valid; the decoded string has at most $10^4$ characters.
\required Time proportional to the length of the output.
\begin{tests}
3[a]2[bc] & aaabcbc\\
3[a2[c]] & accaccacc\\
2[abc]3[cd]ef & abcabccdcdcdef\\
10[a] & aaaaaaaaaa\\
abc & abc\\
\end{tests}

\problem{40}{Remove adjacent duplicates}{Easy}
Repeatedly remove two adjacent equal letters from \cc{s} until no two adjacent letters
are equal; write the result into \cc{out}. The result is unique.
\begin{lstlisting}[style=starter]
void removeAdjacent(char s[], char out[]);
\end{lstlisting}
\permitted Character access; one stack.
\constraints $1\le|s|\le10^5$.
\required $O(n)$ time. (Rescanning the string after each removal is $O(n^2)$.)
\begin{tests}
abbaca & ca\\
azxxzy & ay\\
aa & (empty string)\\
abc & abc\\
\end{tests}

\problem{41}{Sort a stack}{Medium}
Sort the stack \cc{s} so that the smallest item is on top. You may use one additional
stack and scalar variables, nothing else.
\begin{lstlisting}[style=starter]
void sortStack(stack *s);
\end{lstlisting}
\permitted Stack operations only; one extra stack.
\constraints $0\le n\le100$.
\required $O(n^2)$ time.
\begin{tests}
bottom..top 3 1 4 1 5 & popping gives 1 1 3 4 5\\
bottom..top 1 2 3 & popping gives 1 2 3\\
empty & empty\\
\end{tests}

\problem{42}{Validate stack sequences}{Medium}
The distinct numbers \cc{pushed[0..n-1]} are pushed onto an initially empty stack in this
order, and pops may occur at any moment. Return 1 if the pops can produce exactly the
order \cc{popped[0..n-1]}, else 0.
\begin{lstlisting}[style=starter]
int validateStackSequences(int pushed[], int popped[], int n);
\end{lstlisting}
\permitted Array reads; one stack.
\constraints $0\le n\le1000$.
\required $O(n)$ time.
\bonus Prove that the greedy rule ``pop whenever the top is the next wanted item'' never
misses a solution.
\begin{tests}
pushed 1 2 3 4 5, popped 4 5 3 2 1 & 1\\
pushed 1 2 3 4 5, popped 4 3 5 1 2 & 0\\
pushed 2 1 0, popped 1 2 0 & 1\\
\end{tests}

\section{Queues}\label{sec:probqueues}

``Queue operations only'' means \cc{qInit}, \cc{qEnq}, \cc{qDeq}, \cc{qFront},
\cc{qEmpty}, \cc{qSize} of Listing~\ref{lst:queue}: no access to \cc{items[]} or
\cc{front}.

\problem{43}{A queue from two stacks}{Easy}
Implement a FIFO queue using two stacks \cc{in} and \cc{out}: \cc{tsqPush} adds at the
rear, \cc{tsqPop} removes from the front, \cc{tsqPeek} reads the front.
\begin{lstlisting}[style=starter]
typedef struct { stack in; stack out; } twoStackQueue;
void tsqInit(twoStackQueue *q);
void tsqPush(twoStackQueue *q, int x);
int tsqPop(twoStackQueue *q, int *x);    // 0 if empty
int tsqPeek(twoStackQueue *q, int *x);   // 0 if empty
int tsqEmpty(twoStackQueue *q);
\end{lstlisting}
\permitted Stack operations only.
\required Every operation $O(1)$ amortised.
\bonus Prove the amortised bound (Theorem~\ref{thm:twostack}).
\begin{tests}
push 1, push 2, peek, pop, empty & 1, 1, 0\\
then push 3, pop, pop, pop & 2, 3, fails\\
\end{tests}

\problem{44}{A stack from a queue}{Easy}
Implement a LIFO stack using \emph{one} queue.
\begin{lstlisting}[style=starter]
typedef struct { queue q; } queueStack;
void qsInit(queueStack *s);
void qsPush(queueStack *s, int x);
int qsPop(queueStack *s, int *x);    // 0 if empty
int qsTop(queueStack *s, int *x);    // 0 if empty
int qsEmpty(queueStack *s);
\end{lstlisting}
\permitted Queue operations only.
\required One of push and pop in $O(1)$, the other in $O(n)$.
\begin{tests}
push 1, push 2, top, pop, empty & 2, 2, 0\\
then pop, empty & 1, 1\\
push 1, push 2, push 3, pop, push 4, pop, pop & 3, 4, 2\\
\end{tests}

\problem{45}{The Josephus problem}{Medium}
$n$ people stand in a circle, numbered $1..n$. Counting starts at person 1; every
$k$-th person leaves the circle, and counting resumes with the next person. Return the
number of the last person left.
\begin{lstlisting}[style=starter]
int josephus(int n, int k);
\end{lstlisting}
\permitted Queue operations only.
\constraints $1\le n\le99$ (the queue has capacity 100), $1\le k\le1000$.
\required $O(nk)$ time.
\bonus $O(n)$ time and $O(1)$ space without a queue (\cc{josephusFast}), for any $n$.
\begin{tests}
n=7, k=3 & 4\\
n=5, k=2 & 3\\
n=6, k=3 & 1\\
n=41, k=3 & 31\\
n=1, k=5 & 1\\
\end{tests}

\problem{46}{Reverse the first $k$ items of a queue}{Easy}
Reverse the order of the first $k$ items of the queue \cc{q}; the other items keep their
order and stay behind them.
\begin{lstlisting}[style=starter]
void reverseFirstK(queue *q, int k);   // 0 <= k <= size
\end{lstlisting}
\permitted Queue operations on \cc{q}; one stack.
\required $O(n)$ time.
\begin{tests}
{}[1,2,3,4,5], k=3 & [3,2,1,4,5]\\
{}[1,2,3,4,5], k=5 & [5,4,3,2,1]\\
{}[1,2,3,4,5], k=0 & [1,2,3,4,5]\\
{}[10,20,\dots,100], k=5 & [50,40,30,20,10,60,70,80,90,100]\\
\end{tests}

\problem{47}{Number of recent calls}{Easy}
A counter receives pings at times \cc{t} (in milliseconds), each strictly later than the
previous one. \cc{rcPing(r, t)} records the ping and returns how many pings happened in
the inclusive range $[t-3000,\,t]$.
\begin{lstlisting}[style=starter]
typedef struct { queue times; } recentCounter;
void rcInit(recentCounter *r);
int rcPing(recentCounter *r, int t);
\end{lstlisting}
\permitted Queue operations on \cc{times}.
\constraints At most 100 pings fall inside any window.
\required $O(1)$ amortised per ping.
\begin{tests}
ping 1, 100, 3001, 3002 & 1, 2, 3, 3\\
ping 1, 5000, 5001 & 1, 1, 2\\
\end{tests}

\problem{48}{Interleave the two halves of a queue}{Medium}
A queue of even length holds $x_1\dots x_m\,y_1\dots y_m$. Rearrange it into
$x_1y_1x_2y_2\dots x_my_m$.
\begin{lstlisting}[style=starter]
void interleaveHalves(queue *q);
\end{lstlisting}
\permitted Queue operations on \cc{q}; one stack (no second queue, no array).
\required $O(n)$ time.
\begin{tests}
{}[1,2,3,4,5,6,7,8] & [1,5,2,6,3,7,4,8]\\
{}[11,12,\dots,20] & [11,16,12,17,13,18,14,19,15,20]\\
{}[1,2] & [1,2]\\
\end{tests}

\problem{49}{First negative number in every window}{Medium}
For every window \cc{a[i..i+k-1]} of $k$ consecutive items, write into \cc{out} the
first negative number of the window, or 0 if it has none. Return the number of windows,
$n-k+1$.
\begin{lstlisting}[style=starter]
int firstNegatives(int a[], int n, int k, int out[]);
\end{lstlisting}
\permitted Array access to \cc{a} and \cc{out}; one queue.
\constraints $1\le k\le n\le10^5$.
\required $O(n)$ time ($O(nk)$ earns partial credit).
\begin{tests}
{}[12,-1,-7,8,-15,30,16,28], k=3 & [-1,-1,-7,-15,-15,0]\\
{}[-8,2,3,-6,10], k=2 & [-8,0,-6,-6]\\
{}[1,2,3], k=3 & [0]\\
\end{tests}

\section{Deques}\label{sec:probdeques}

``Deque operations only'' means \cc{push\_front}, \cc{push\_back}, \cc{pop\_front},
\cc{pop\_back}, \cc{front}, \cc{back}, \cc{empty} and \cc{size} of \cc{deque<int>}.

\problem{50}{Design a circular deque}{Medium}
Implement a deque of capacity $k$ in a circular array. The insert and delete operations
return \cc{true} on success and \cc{false} if the deque is full (insert) or empty
(delete); \cc{getFront} and \cc{getRear} return $-1$ if the deque is empty.
\begin{lstlisting}[style=starter]
class CircularDeque {
public:
    int items[1000];
    int cap, first, count;
    CircularDeque(int k);       // called as  CircularDeque d(3);
    bool insertFront(int x);
    bool insertLast(int x);
    bool deleteFront();
    bool deleteLast();
    int getFront();
    int getRear();
    bool isEmpty();
    bool isFull();
};
\end{lstlisting}
\permitted Array access to the fields.
\constraints $1\le k\le1000$.
\required Every operation $O(1)$; no shifting.
\begin{tests}
k=3: insertLast 1, insertLast 2, insertFront 3, insertFront 4 & true, true, true, false\\
getRear, isFull, deleteLast, insertFront 4, getFront & 2, true, true, true, 4\\
k=1: getFront, insertFront 7, insertLast 8, getRear & -1, true, false, 7\\
\end{tests}

\problem{51}{A linked deque with a sentinel}{Medium}
Implement a deque on a circular doubly linked list with one sentinel node:
\cc{sentinel->next} is the first item and \cc{sentinel->prev} the last; when the deque is
empty both point to the sentinel. Write the two helpers first and implement the
operations with them, without any \cc{if} that tests for \cc{NULL}.
\begin{lstlisting}[style=starter]
class dqnode {
public:
    int item;
    dqnode *prev, *next;
    dqnode(int x);
};
class LinkedDeque {
public:
    dqnode *sentinel;
    int count;
    LinkedDeque();                          // creates the sentinel
    void linkBetween(dqnode *a, dqnode *b, int x);
    void unlink(dqnode *n);
    void push_front(int x);   void push_back(int x);
    void pop_front();         void pop_back();
    int front();              int back();
    bool empty();             int size();
    void clear();                           // deletes every item node
};
\end{lstlisting}
\permitted \cc{new}, \cc{delete}, pointer manipulation.
\required Every operation $O(1)$ (\cc{clear}: $O(n)$); no leaks.
\bonus Explain what goes wrong after \cc{LinkedDeque e = d;}.
\begin{tests}
push\_back 1, push\_back 2, push\_front 0 & 0 1 2\\
back, pop\_back; front, pop\_front & 2, 0\\
front, back & 1, 1\\
pop\_back; empty(); pop\_back & (removed); true; nothing happens\\
\end{tests}

\problem{52}{Valid palindrome}{Easy}
Return 1 if the string \cc{s} reads the same forwards and backwards when only its letters
and digits are kept and upper case is treated as lower case; else 0.
\begin{lstlisting}[style=starter]
int isPalindromeText(char s[]);
\end{lstlisting}
\permitted Character reads of \cc{s}; deque operations. No \cc{ctype.h} (write the tests
yourself).
\constraints $1\le|s|\le100$ letters and digits after filtering (the deque has capacity
100).
\required $O(n)$ time.
\bonus $O(1)$ extra space with two indices instead of a deque.
\begin{tests}
A man, a plan, a canal: Panama & 1\\
race a car & 0\\
{" "} (one space) & 1\\
0P & 0\\
Was it a car or a cat I saw? & 1\\
\end{tests}

\problem{53}{Sliding window maximum}{Medium}
For every window \cc{a[i..i+k-1]} write its maximum into \cc{out}; return the number of
windows, $n-k+1$.
\begin{lstlisting}[style=starter]
int maxSlidingWindow(int a[], int n, int k, int out[]);
\end{lstlisting}
\permitted Array access to \cc{a} and \cc{out}; one deque.
\constraints $1\le k\le n\le10^5$, $k\le100$.
\required $O(n)$ time ($O(nk)$ earns partial credit).
\begin{tests}
{}[1,3,-1,-3,5,3,6,7], k=3 & [3,3,5,5,6,7]\\
{}[1], k=1 & [1]\\
{}[9,1,8,2,7,3], k=2 & [9,8,8,7,7]\\
{}[4,2,12,3,8,8,1,5], k=3 & [12,12,12,8,8,8]\\
\end{tests}

\problem{54}{Longest subarray with absolute difference at most a limit}{Medium}
Return the length of the longest non-empty contiguous subarray in which the difference
between the largest and the smallest item is at most \cc{limit}.
\begin{lstlisting}[style=starter]
int longestSubarray(int a[], int n, int limit);
\end{lstlisting}
\permitted Array reads; two deques.
\constraints $1\le n\le10^5$, $0\le\mathit{limit}$; no valid window is longer than 100.
\required $O(n)$ time.
\begin{tests}
{}[8,2,4,7], limit 4 & 2\\
{}[10,1,2,4,7,2], limit 5 & 4 ([2,4,7,2])\\
{}[4,2,2,2,4,4,2,2], limit 0 & 3\\
\end{tests}

\problem{55}{Front-middle-back queue}{Medium}
Design a queue that supports \cc{push} and \cc{pop} at the front, in the middle and at
the back. When there are two middle positions, the \emph{front-most} one is used:
\cc{pushMiddle} inserts so that the new item becomes item $\lfloor n/2\rfloor$ (0-based,
$n$ the size before), and \cc{popMiddle} removes item $\lfloor(n-1)/2\rfloor$. The pops
return the removed item, or $-1$ if the queue is empty.
\begin{lstlisting}[style=starter]
class FrontMiddleBackQueue {
public:
    deque<int> left;
    deque<int> right;
    void pushFront(int x);
    void pushMiddle(int x);
    void pushBack(int x);
    int popFront();
    int popMiddle();
    int popBack();
};
\end{lstlisting}
\permitted Deque operations on the two member deques.
\required Every operation $O(1)$.
\begin{tests}
pushFront 1, pushBack 2, pushMiddle 3, pushMiddle 4 & queue [1,4,3,2]\\
popFront, popMiddle, popMiddle, popBack, popFront & 1, 3, 4, 2, -1\\
\end{tests}

\problem{56}{ArrayDeque: positional insertion and removal}{Medium}
In the ArrayDeque of the course notes, item $i$ ($0\le i<n$) lives at
\cc{a[(j+i) \% 100]}. Implement \cc{add(i, x)}, which inserts $x$ so that it becomes
item $i$, and \cc{remove(i)}, which removes item $i$ and returns it. Each must move the
\emph{shorter} side of the deque: the items before position $i$ (moving \cc{j}) or the
items after it.
\begin{lstlisting}[style=starter]
class ArrayDeque {
public:
    int a[100];
    int j;
    int n;
    ArrayDeque();                // j = 0, n = 0
    int size();
    int get(int i);
    bool add(int i, int x);      // false on failure
    int remove(int i);           // returns the removed item, -1 on failure
};
\end{lstlisting}
\permitted Array access to the fields.
\required $O(1+\min(i,n-i))$ time per operation.
\begin{tests}
add(0,1), add(1,2), add(2,3), add(0,0) & [0,1,2,3]\\
add(2,9) & [0,1,9,2,3]\\
remove(1); get(2) & returns 1, [0,9,2,3]; 2\\
add(5,7); remove(3) & false; returns 3, [0,9,2]\\
\end{tests}

\problem{57}{Jump game VI}{Medium}
You start on index 0 of the array \cc{a} and must reach index $n-1$. From index $i$ you
may jump to any index $i+1,\dots,i+k$ (inside the array). Your score is the sum of
\cc{a[j]} over all indices $j$ you land on, including the first and the last. Return the
maximum score.
\begin{lstlisting}[style=starter]
int maxResult(int a[], int n, int k);
\end{lstlisting}
\permitted Array access; one deque.
\constraints $1\le n\le10^4$, $1\le k\le98$, $|a[i]|\le10^4$.
\required $O(n)$ time ($O(nk)$ earns partial credit).
\begin{tests}
{}[1,-1,-2,4,-7,3], k=2 & 7 (path 1, -1, 4, 3)\\
{}[10,-5,-2,4,0,3], k=3 & 17\\
{}[1,-5,-20,4,-1,3,-6,-3], k=2 & 0\\
\end{tests}

